% GENERATED FILE. Edit canonical lessons, then run npm run generate:books.
% canonical-source-hash: fnv1a64:dcbb1f791641a487
% canonical-lessons: ZH-C375-qiong, ZH-C375-shehui, ZH-W375-she-society, ZH-C375-juan, ZH-C375-zhiyuanzhe, ZH-C375-huxiang

\chapter{Helping Others}
\label{ch:zh-helping-others}
\hlchaptermodality{375}

\begin{chapteropening}
\textbf{By the end of this chapter:} \emph{I can talk about being poor, society, giving money, volunteers and helping each other.}

Complete the last lesson of chapter 375: I can talk about being poor, society, giving money, volunteers and helping each other.
\end{chapteropening}

\section[qióng]{qióng — poor (having little money)}

\label{lesson:ZH-C375-qiong}



\begin{quote}
Before the new one: say the Mandarin for an ethnic group, then the Mandarin for to develop and grow.
\end{quote}

\subsection*{You'll want to know: qióng}
\emph{qióng} — \textquotedblleft{}poor (having little money)\textquotedblright{}.

\begin{quote}
\emph{yǐqián zhèr de rén dōu hěn qióng} — People here used to be poor.
\end{quote}

Its characters come later; for now it is a word for the ear.

\begin{cousinweb}
\emph{qióng} is inherited from Middle Chinese \emph{gjuwng} (Guangyun, in Baxter's transcription), glossed 'to come to the very end'. How running out came to mean being poor is not given there.
\end{cousinweb}

\begin{grammarlens}[title={What you've built}]
One more word for giving reasons.
\end{grammarlens}

\subsection*{Your turn}
\begin{itemize}
\raggedright
  \item \emph{Say it:} \emph{qióng}
  \item \emph{Say it:} \emph{qióng}, once more
  \item \emph{Say it:} \emph{yǐqián zhèr de rén dōu hěn qióng} — People here used to be poor.
  \item \emph{Recall:} say the Mandarin for an ethnic group, then the Mandarin for to develop and grow, then say \emph{qióng} again
\end{itemize}

\begin{tcolorbox}[breakable,colback=teal!4,colframe=teal!35!black,title={Before you move on}]
What does qióng mean? (\textquotedblleft{}Poor (having little money)\textquotedblright{}.) Say it once more.
\end{tcolorbox}
\section[\texorpdfstring{\zh{社会} (shèhuì)}{社会 (shèhuì)}]{\zh{社会} (shèhuì) — society}

\label{lesson:ZH-C375-shehui}



\begin{quote}
Before the new one: say the Mandarin for to develop and grow, then the Mandarin for poor (having little money).
\end{quote}

\subsection*{You'll want to know: \zh{社会}}
\textbf{\zh{社会}} — \emph{shèhuì} — \textquotedblleft{}society\textquotedblright{}.

\begin{quote}
\emph{wǒmen de shèhuì xūyào rèxīn de rén} — Our society needs warm-hearted people.
\end{quote}

The new character, in pieces: \zh{社} — seven strokes in two pieces: on the left the sign piece, four strokes; on the right a short bar, a vertical and a long bar, the earth piece at the foot of \zh{在}. The next lesson puts it in your hand.

\begin{cousinweb}
\textbf{\zh{社会}} (\emph{shèhuì}) is built from \zh{社} 'the altar of the soil' and \zh{会} 'to meet', both inherited from Middle Chinese, \emph{dzyaeX} and \emph{hwajH} (Guangyun, in Baxter's transcription), which lists the first with the god of grain and glosses the second 'to come together'. People meeting at the village altar is usually given as the start; the sense 'society' is not given there.
\end{cousinweb}

\begin{grammarlens}[title={What you've built}]
One more word for giving reasons.
\end{grammarlens}

\subsection*{Your turn}
\begin{itemize}
\raggedright
  \item \emph{Say it:} \emph{shèhuì}
  \item \emph{Say it:} \emph{shèhuì}, once more
  \item \emph{Say it:} \emph{wǒmen de shèhuì xūyào rèxīn de rén} — Our society needs warm-hearted people.
  \item \emph{Recall:} say the Mandarin for to develop and grow, then the Mandarin for poor (having little money), then say \emph{shèhuì} again
\end{itemize}

\begin{tcolorbox}[breakable,colback=teal!4,colframe=teal!35!black,title={Before you move on}]
What does \zh{社会} mean? (\textquotedblleft{}Society\textquotedblright{}.) Say it once more.
\end{tcolorbox}
\section[\texorpdfstring{\zh{社} (shè)}{社 (shè)}]{\zh{社} — the sign piece, then the earth piece}

\label{lesson:ZH-W375-she-society}



\begin{quote}
The word first, then the character inside it.

\begin{itemize}
\raggedright
  \item \emph{Recall:} say \emph{shèhuì} — \textbf{R1}, one lesson back, before it decays
\end{itemize}

\begin{quote}
\zh{在} \zh{会}
\end{quote}

\emph{Read it:} each one aloud, with nothing said first — \emph{zài}, \emph{huì}
\end{quote}

\subsection*{Script — the sign piece, then the earth piece}
\hlblockfigure{\detokenize{figures/ZH-W375-she-society-filmstrip.pdf}}{How \zh{社} is written, stroke by stroke}

\begin{quote}
\zh{社}
\end{quote}

\textbf{Two pieces, left before right.} Count the pieces, not the strokes.

\textbf{The left piece, the sign piece, four strokes.} A short dot at the top. A bar that runs right, turns and falls to the lower left, all one stroke. A vertical down from under it. Then a short dot to the right of the vertical.

\textbf{The right piece, the earth piece, three strokes,} as at the foot of \zh{在}: a short bar, a vertical down through it, then a long bar across the foot.

4 + 3 = seven.

\subsection*{Your turn}
\begin{itemize}
\raggedright
  \item \emph{Write it:} \zh{社}, all seven, with the model in front of you
  \item \emph{Cover:} write it again from memory, and count the pieces, not the strokes
\end{itemize}

\begin{tcolorbox}[breakable,colback=teal!4,colframe=teal!35!black,title={Before you move on}]
How many pieces in \zh{社}? (\textbf{Two}.) How many strokes in the left piece? (\textbf{Four}.) Where have you made the right piece before? (\textbf{At the foot of \zh{在}}.)

Source: \href{https://raw.githubusercontent.com/chanind/hanzi-writer-data/68d10a4b21150cae5e1ebbd223eed289cf32d90c/data/\%E7\%A4\%BE.json}{Hanzi Writer Data \zh{社}, pinned stroke paths and medians}.
\end{tcolorbox}
\section[juān]{juān — to donate}

\label{lesson:ZH-C375-juan}



\begin{quote}
Before the new one: say the Mandarin for poor (having little money), then the Mandarin for society.
\end{quote}

\subsection*{You'll want to know: juān}
\emph{juān} — \textquotedblleft{}to donate\textquotedblright{}.

\begin{quote}
\emph{tā měi nián dōu juān yìxiē qián} — He gives some money away every year.
\end{quote}

Its characters come later; for now it is a word for the ear.

\begin{cousinweb}
\emph{juān} is inherited from Middle Chinese \emph{ywen} (Guangyun, in Baxter's transcription), glossed 'to throw away'. Giving up money for others is the later sense, and it is not given there.
\end{cousinweb}

\begin{grammarlens}[title={What you've built}]
One more word for giving reasons.
\end{grammarlens}

\subsection*{Your turn}
\begin{itemize}
\raggedright
  \item \emph{Say it:} \emph{juān}
  \item \emph{Say it:} \emph{juān}, once more
  \item \emph{Say it:} \emph{tā měi nián dōu juān yìxiē qián} — He gives some money away every year.
  \item \emph{Recall:} say the Mandarin for poor (having little money), then the Mandarin for society, then say \emph{juān} again
\end{itemize}

\begin{tcolorbox}[breakable,colback=teal!4,colframe=teal!35!black,title={Before you move on}]
What does juān mean? (\textquotedblleft{}To donate\textquotedblright{}.) Say it once more.
\end{tcolorbox}
\section[zhìyuànzhě]{zhìyuànzhě — volunteer}

\label{lesson:ZH-C375-zhiyuanzhe}



\begin{quote}
Before the new one: say the Mandarin for society, then the Mandarin for to donate (money).
\end{quote}

\subsection*{You'll want to know: zhìyuànzhě}
\emph{zhìyuànzhě} — \textquotedblleft{}volunteer\textquotedblright{}.

\begin{quote}
\emph{wǒ xiǎng dāng zhìyuànzhě yīnwèi wǒ xiǎng bāng biérén} — I want to be a volunteer because I want to help other people.
\end{quote}

Its characters come later; for now it is a word for the ear.

\begin{cousinweb}
\emph{zhìyuànzhě} is built from \emph{zhì} 'what the heart longs for', \emph{yuàn} 'to wish' and \emph{zhě} 'one who', all inherited from Middle Chinese, \emph{tsyiH}, \emph{ngjwonH} and \emph{tsyaeX} (Guangyun, in Baxter's transcription), which lists the last as a grammar word. A volunteer is one who serves of their own will.
\end{cousinweb}

\begin{grammarlens}[title={What you've built}]
One more word for giving reasons.
\end{grammarlens}

\subsection*{Your turn}
\begin{itemize}
\raggedright
  \item \emph{Say it:} \emph{zhìyuànzhě}
  \item \emph{Say it:} \emph{zhìyuànzhě}, once more
  \item \emph{Say it:} \emph{wǒ xiǎng dāng zhìyuànzhě yīnwèi wǒ xiǎng bāng biérén} — I want to be a volunteer because I want to help other people.
  \item \emph{Recall:} say the Mandarin for society, then the Mandarin for to donate (money), then say \emph{zhìyuànzhě} again
\end{itemize}

\begin{tcolorbox}[breakable,colback=teal!4,colframe=teal!35!black,title={Before you move on}]
What does zhìyuànzhě mean? (\textquotedblleft{}Volunteer\textquotedblright{}.) Say it once more.
\end{tcolorbox}
\section[hùxiāng]{hùxiāng — each other; mutually}

\label{lesson:ZH-C375-huxiang}



\begin{quote}
Before the new one: say the Mandarin for to donate (money), then the Mandarin for a volunteer.
\end{quote}

\subsection*{You'll want to know: hùxiāng}
\emph{hùxiāng} — \textquotedblleft{}each other; mutually\textquotedblright{}.

\begin{quote}
\emph{dàjiā yīnggāi hùxiāng bāngmáng} — Everyone should help each other.
\end{quote}

Its characters come later; for now it is a word for the ear.

\begin{cousinweb}
\emph{hùxiāng} is built from \emph{hù} and \emph{xiāng}, both inherited from Middle Chinese, \emph{huH} and \emph{sjang} (Guangyun, in Baxter's transcription), which glosses the first 'interlocking' and the second 'together'.
\end{cousinweb}

\begin{grammarlens}[title={What you've built}]
One more word for giving reasons.
\end{grammarlens}

\subsection*{Your turn}
\begin{itemize}
\raggedright
  \item \emph{Say it:} \emph{hùxiāng}
  \item \emph{Say it:} \emph{hùxiāng}, once more
  \item \emph{Say it:} \emph{dàjiā yīnggāi hùxiāng bāngmáng} — Everyone should help each other.
  \item \emph{Recall:} say the Mandarin for to donate (money), then the Mandarin for a volunteer, then say \emph{hùxiāng} again
\end{itemize}

\begin{tcolorbox}[breakable,colback=teal!4,colframe=teal!35!black,title={Before you move on}]
What does hùxiāng mean? (\textquotedblleft{}Each other; mutually\textquotedblright{}.) Say it once more.
\end{tcolorbox}
